% Unit 130 terjemahan Bahasa Indonesia dari chapter9.tex baris 1457--1625.
% Sumber: Methods of Algebra, Volume 2, commit
% 9a5803ff77dd3257484cb177f851a73770a59dd3 (CC BY 4.0).
% stable-unit-id: o014.aljabr2.chapter9.hopf-reconstruction
% Cakupan: Bagian 9.6, rekonstruksi koaljabar, bialjabar, dan antipoda.

\section{Rekonstruksi Aljabar Hopf}\label{sec:Hopf-reconstruction}
Bagian ini melanjutkan pembahasan dalam \S\ref{sec:Tannaka-coend}. Teorema
\ref{prop:Tannaka-aux1} membangun koaljabar $\coEnd(\omega)$ dari data
$(\mathcal{A}, \omega)$ dan mengidentifikasikan $\mathcal{A}$ dengan
$\catesubd{Comod}{\mathrm{f}}\coEnd(\omega)$. Situasi paling mendasar ialah
ketika $B = \Bbbk$ merupakan medan. Untuk suatu koaljabar-$\Bbbk$ $L$, ambil
\begin{gather*}
	\mathcal{A} := \catesubd{Comod}{\mathrm{f}}L , \\
	\omega := U: \catesubd{Comod}{\mathrm{f}}L \xrightarrow{\text{lupa}} \cate{Vect}_{\mathrm{f}}(\Bbbk).
\end{gather*}

Karena $\Bbbk$ medan, Proposisi \ref{prop:Comod-abelian} dengan $B = \Bbbk$
menunjukkan bahwa $\cated{Comod}L$ merupakan kategori abelian, sedangkan
$\catesubd{Comod}{\mathrm{f}}L$ yang dibentuk oleh komodul berdimensi hingga
jelas merupakan subkategori abelian hingga secara lokal. Kini wajar untuk bertanya:
apakah $\coEnd(\omega)$ yang bersesuaian merekonstruksi data semula $L$, atau
justru menghasilkan koaljabar baru? Jawabannya, sebagaimana diduga, ialah
``rekonstruksi'', tetapi diperlukan sedikit argumen.

Langkah pertama ialah memperjelas cara membandingkan $\coEnd(\omega)$ dan
$L$. Karena $L$ berkoaksi dari kanan pada setiap $\omega(X)$, sifat universal
menginduksi homomorfisme koaljabar
\begin{equation}\label{eqn:u-comparison}
	u: \coEnd(\omega) \to L.
\end{equation}
Persoalan semula sama dengan menanyakan apakah $u$ merupakan isomorfisme.
Kita terlebih dahulu membuktikan sebuah hasil sederhana.

\begin{lemma}\label{prop:comodule-fd-union}
	Misalkan $\Bbbk$ medan dan $L$ koaljabar-$\Bbbk$, dengan homomorfisme
	komultiplikasi dinotasikan oleh $\Delta$ dan homomorfisme kounit oleh
	$\epsilon$. Setiap komodul-$L$ kanan $V$ merupakan gabungan dari
	subkomodul-subkomodul berdimensi hingga.
\end{lemma}
\begin{proof}
	Misalkan struktur komodul diberikan oleh
	$\rho: V \to V \dotimes{\Bbbk} L$. Untuk setiap
	$x \in V \dotimes{\Bbbk} L$, definisikan subruang
	\[ V(x) := \lrangle{(\identity \otimes \lambda)(x) : \lambda \in \Hom_{\Bbbk}(L, \Bbbk) } \; \subset V ; \]
	jika $x$ dituliskan secara konkret sebagai
	$\sum_{i=1}^n v_i \otimes \ell_i$, dengan
	$\ell_1, \ldots, \ell_n$ bebas linear, maka jelas bahwa
	$V(x) = \lrangle{v_1, \ldots, v_n}$, sehingga
	$x \in V(x) \dotimes{\Bbbk} L$.
	
	Ambil $v \in V$. Dengan mengambil $x = \rho(v)$ dan
	$\lambda = \epsilon$, diperoleh $v \in V(\rho(v))$. Selanjutnya kita
	tunjukkan bahwa
	$\rho(V(\rho(v))) \subset V(\rho(v)) \dotimes{\Bbbk} L$. Tuliskan lagi
	$\rho(v) = \sum_i v_i \otimes \ell_i$; persamaan
	$(\rho \otimes \identity)\rho(v) = (\identity \otimes \Delta)\rho(v)$
	dalam definisi komodul dapat ditulis sebagai persamaan dalam
	$V \dotimes{\Bbbk} L \dotimes{\Bbbk} L$
	\[ \sum_i \rho(v_i) \otimes \ell_i = \sum_i v_i \otimes \Delta(\ell_i). \]
	
	Perluas $\ell_1, \ldots, \ell_n$ menjadi suatu basis, lalu kontraksikan
	dengan basis dualnya. Diperoleh
	$\rho(v_i) \in V(\rho(v)) \dotimes{\Bbbk} L$ untuk
	$i = 1, \ldots, n$.
	
	Dengan demikian, setiap subruang berdimensi hingga
	$\lrangle{v_1, \ldots, v_m}$ dari $V$ termuat dalam subkomodul
	$\sum_{i=1}^m V(\rho(v_i))$, yang jelas berdimensi hingga.
\end{proof}

\begin{lemma}\label{prop:reconstruction-cat-aux}
	Dengan notasi di atas, $u: \coEnd(\omega) \to L$ dalam
	\eqref{eqn:u-comparison} merupakan isomorfisme koaljabar.
\end{lemma}
\begin{proof}
	Kita tunjukkan bahwa semua komodul-$L$ kanan $V$, beserta semua
	homomorfisme antarkomodul, secara alami terangkat menjadi
	komodul-$\coEnd(\omega)$ kanan. Untuk $\dim_{\Bbbk} V < \infty$, hal ini
	merupakan definisi---penerapan Definisi--Proposisi
	\ref{def:cogebra-L} pada
	$\mathcal{A} = \catesubd{Comod}{\mathrm{f}}L$ dan $\omega = U$;
	kasus umum kemudian mengikuti dari Lema
	\ref{prop:comodule-fd-union}.
	
	Sekarang ambil $V=L$ dengan struktur komodul-$L$ kanan yang ditentukan
	oleh $\Delta$. Struktur komodul-$\coEnd(\omega)$ hasil pengangkatan
	bersesuaian dengan
	$\tilde{\Delta}: L \to L \dotimes{\Bbbk} \coEnd(\omega)$. Ini sama
	dengan mengatakan bahwa bagian segitiga dari diagram berikut komutatif:
	\[\begin{tikzcd}
		L \arrow[r, "{\tilde{\Delta}}"] \arrow[rd, "\Delta"'] & L \dotimes{\Bbbk} \coEnd(\omega) \arrow[r, "{\epsilon \otimes \identity}"] \arrow[d, "{\identity \otimes u}"] & \coEnd(\omega) \arrow[d, "u"] \\
		& L \dotimes{\Bbbk} L \arrow[r, "{\epsilon \otimes \identity}"'] & L,
	\end{tikzcd}\]
	dan bagian perseginya jelas juga komutatif. Nyatakan komposisi pada baris
	pertama sebagai $\gamma$; maka $u\gamma = \identity_L$. Cukup dibuktikan
	bahwa $\gamma$ surjektif.
	
	Misalkan $\rho: V \to V \dotimes{\Bbbk} L$ suatu komodul-$L$ kanan
	berdimensi hingga. Diagram komutatif dalam definisi
	\[\begin{tikzcd}
		V \arrow[r, "\rho"] \arrow[d, "\rho"'] & V \dotimes{\Bbbk} L \arrow[d, "{\identity \otimes \Delta}"] \\
		V \dotimes{\Bbbk} L \arrow[r, "{\rho \otimes \identity}"'] & V \dotimes{\Bbbk} L \dotimes{\Bbbk} L
	\end{tikzcd}\]
	juga sama dengan mengatakan bahwa $\rho$ merupakan homomorfisme
	komodul-$L$, asalkan $V \dotimes{\Bbbk} L$ dibekali struktur komodul
	yang berasal dari $L$ dan $\Delta$. Sekarang angkat $V$ secara alami
	menjadi komodul-$\coEnd(\omega)$, yang bersesuaian dengan
	$\tilde{\rho}: V \to V \dotimes{\Bbbk} \coEnd(\omega)$; jika komodul
	$V \dotimes{\Bbbk} L$ diangkat dengan cara serupa, struktur itu
	bersesuaian dengan
	$\identity \otimes \tilde{\Delta}: V \dotimes{\Bbbk} L \to V \dotimes{\Bbbk} L \dotimes{\Bbbk} \coEnd(\omega)$. Homomorfisme komodul
	$\rho$ juga terangkat ke tingkat $\coEnd(\omega)$, yang sama dengan
	mengatakan bahwa persegi dalam diagram berikut komutatif:
	\[\begin{tikzcd}
		V \arrow[r, "\rho"] \arrow[d, "{\tilde{\rho}}"'] & V \dotimes{\Bbbk} L \arrow[d, "{\identity \otimes \tilde{\Delta}}"] & \\
		V \dotimes{\Bbbk} \coEnd(\omega) \arrow[r, "{\rho \otimes \identity}"' inner sep=0.8em] & V \dotimes{\Bbbk} L \dotimes{\Bbbk} \coEnd(\omega) \arrow[r, "{\identity \otimes \epsilon \otimes \identity}"' inner sep=0.8em] & V \dotimes{\Bbbk} \coEnd(\omega)
	\end{tikzcd}\]
	Karena komposisi baris kedua adalah $\identity$, diperoleh
	$\tilde{\rho} = (\identity_V \otimes \gamma)\rho$. Dapat diperiksa
	bahwa pemetaan
	$\alpha_V: V^\vee \dotimes{\Bbbk} V \to \coEnd(\omega)$ yang diinduksi
	oleh $\tilde{\rho}$ dengan demikian berfaktor sebagai
	\[ V^\vee \dotimes{\Bbbk} V \xrightarrow{\text{diinduksi oleh $\rho$}} L \xrightarrow{\gamma} \coEnd(\omega). \]
	
	Dengan mengingat kembali konstruksi $\coEnd(\omega)$, ketika $V$
	bervariasi, $\Image(\alpha_V)$ membangkitkan seluruh $\coEnd(\omega)$.
	Jadi $\gamma$ surjektif.
\end{proof}

Dapat diduga bahwa jika $L$ bialjabar, maka \eqref{eqn:u-comparison} juga
merupakan isomorfisme bialjabar. Argumen yang diperlukan seluruhnya rutin
dan tidak perlu dijabarkan.

\begin{definition}\label{def:Fib-cat}
	\index[sym1]{Fibk@$\cate{Fib}(\Bbbk)$, $\cate{Fib}^{\otimes}(\Bbbk)$}
	Untuk medan $\Bbbk$ yang telah dipilih, definisikan kategori
	$\cate{Fib}(\Bbbk)$ sebagai berikut.
	\begin{itemize}
		\item Objeknya ialah data $(\mathcal{A}, \omega)$, dengan
		$\mathcal{A}$ kategori abelian hingga secara lokal (sehingga $\mathcal{A}$ kecil
		secara esensial), dan
		$\omega: \mathcal{A} \to \cate{Vect}_{\mathrm{f}}(\Bbbk)$ funktor
		setia eksak.
		\item Morfisme dari $(\mathcal{A}, \omega)$ ke
		$(\mathcal{A}', \omega')$ ialah data $(F, \alpha)$, dengan
		$F: \mathcal{A} \to \mathcal{A}'$ funktor dan
		$\alpha: \omega \rightiso \omega' F$ isomorfisme antarfunktor.
		Identitas dan komposisi didefinisikan dengan cara yang jelas.
	\end{itemize}
	Dengan cara serupa, definisikan kategori
	$\cate{Fib}^{\otimes}(\Bbbk)$, dengan syarat tambahan berikut:
	\begin{itemize}
		\item $\mathcal{A}$ dalam objek $(\mathcal{A}, \omega)$ harus
		dilengkapi struktur monoidal, dan $\omega$ harus merupakan funktor
		monoidal;
		\item $F$ dalam suatu morfisme juga harus merupakan funktor monoidal,
		dan $\alpha$ harus merupakan isomorfisme antarfunktor monoidal.
	\end{itemize}
\end{definition}

Perhatikan bahwa jika $(\mathcal{A}, \omega)$ objek
$\cate{Fib}^{\otimes}(\Bbbk)$, syarat setia dan eksak menyiratkan bahwa
$\End_{\mathcal{A}}(\munit) \to \Bbbk$ merupakan monomorfisme
aljabar-$\Bbbk$; jadi dalam hal ini
$\End_{\mathcal{A}}(\munit) \rightiso \Bbbk$.

Di sisi lain, nyatakan kategori koaljabar-$\Bbbk$ (atau bialjabar-$\Bbbk$)
beserta homomorfisme di antaranya sebagai $\Bbbk\dcate{coAlg}$ (atau
$\Bbbk\dcate{biAlg}$).
\begin{itemize}
	\item Terdapat funktor yang jelas
	$\cate{Fib}(\Bbbk) \to \Bbbk\dcate{coAlg}$ dan
	$\cate{Fib}^{\otimes}(\Bbbk) \to \Bbbk\dcate{biAlg}$. Keduanya
	memetakan objek $(\mathcal{A}, \omega)$ ke $\coEnd(\omega)$, sedangkan
	morfisme $(\mathcal{A}, \omega) \to (\mathcal{A}', \omega')$
	menginduksi homomorfisme koaljabar atau bialjabar
	$\coEnd(\omega) \simeq \coEnd(\omega' F) \to \coEnd(\omega')$.
	\item Terdapat pula funktor ke arah sebaliknya: untuk koaljabar atau
	bialjabar $L$, ambil data yang bersesuaian
	$\left(\catesubd{Comod}{\mathrm{f}}L, U \right)$; pada tingkat morfisme,
	untuk homomorfisme $\phi: L \to L'$, komodul-$L$ kanan $(M, \rho)$
	menjadi komodul-$L'$ kanan melalui
	$M \xrightarrow{(\identity \otimes \phi) \rho} M \otimes L'$.
	Ini memberikan
	$\left(\catesubd{Comod}{\mathrm{f}}L, U \right) \to \left(\catesubd{Comod}{\mathrm{f}}L', U' \right)$.
\end{itemize}

\begin{theorem}\label{prop:Tannaka-reconstruction}
	\index{teorema rekonstruksi}
	Misalkan $\Bbbk$ medan. Konstruksi di atas memberikan dua pasang funktor
	yang saling menjadi kuasi-invers
	\[\begin{tikzcd}[row sep=tiny]
		\cate{Fib}(\Bbbk) \arrow[shift left, r] & \Bbbk\dcate{coAlg}, \arrow[shift left, l] \\
		\cate{Fib}^{\otimes}(\Bbbk) \arrow[shift left, r] & \Bbbk\dcate{biAlg}. \arrow[shift left, l]
	\end{tikzcd}\]
\end{theorem}
\begin{proof}
	Tuliskan funktor-funktor dalam kedua arah sebagai
	$Z: \star \leftrightarrows \star :Y$. Untuk pasangan pertama, isomorfisme
	$u$ dalam Lema \ref{prop:reconstruction-cat-aux} memberikan
	$ZY \rightiso \identity$, sedangkan ekuivalensi $\overline{\omega}$
	dalam Teorema \ref{prop:Tannaka-aux1} (i) memberikan
	$\identity \rightiso YZ$.
	
	Untuk pasangan kedua, karena telah diketahui bahwa ketika $L$ bialjabar,
	$u$ dalam \eqref{eqn:u-comparison} juga merupakan isomorfisme bialjabar,
	dalam situasi monoidal masih berlaku $ZY \rightiso \identity$; sementara
	itu, $\identity \rightiso YZ$ kini merupakan isi Teorema
	\ref{prop:Tannaka-aux1} (ii).
\end{proof}

Sekarang kita tambahkan dualitas yang diperkenalkan dalam
\S\ref{sec:monoidal-dual} ke dalam kerangka Teorema Rekonstruksi
\ref{prop:Tannaka-aux1}. Seperti biasa, kita hanya meninjau situasi sederhana
ketika $B = \Bbbk$ merupakan medan.

\begin{theorem}\label{prop:Tannaka-aux2}
	Di bawah asumsi Teorema \ref{prop:Tannaka-reconstruction}, andaikan lebih
	lanjut bahwa $(\mathcal{A}, \omega)$ merupakan objek
	$\cate{Fib}^{\otimes}(\Bbbk)$, serta $\mathcal{A}$ kaku kiri dan kaku
	kanan. Maka $\coEnd(\omega)$ merupakan aljabar Hopf.
\end{theorem}
\begin{proof}
	Cukup ditunjukkan bahwa bialjabar $\coEnd(\omega)$ memiliki antipoda.
	Gunakan kekakuan kiri $\mathcal{A}$ untuk mendefinisikan funktor
	$X \mapsto {}^* X$. Untuk setiap $X \in \Obj(\mathcal{A})$, karena
	funktor monoidal mempertahankan dual, diperoleh
	$\omega({}^* X) \simeq \omega(X)^\vee$; karena struktur monoidal
	$\cate{Vect}_{\mathrm{f}}(\Bbbk)$ simetris, juga diperoleh
	$\omega(X) \simeq \omega({}^* X)^\vee$. Sekarang tinjau
	\[ \phi \in \omega(X)^\vee \simeq \omega({}^* X), \quad x \in \omega(X) \simeq \omega({}^* X)^\vee , \]
	dengan $X \in \Obj(\mathcal{A})$ sembarang. Dalam notasi Konvensi
	\ref{con:coend-bracket}, baik $[\phi \otimes x]$ maupun
	$[x \otimes \phi]$ dengan demikian merupakan unsur $\coEnd(\omega)$.
	Kita tunjukkan bahwa
	\begin{equation}\label{eqn:Tannaka-aux2}
		\text{terdapat pemetaan linear}\; S: \coEnd(\omega) \to \coEnd(\omega), \;\text{yang ditentukan oleh}\; S[\phi \otimes x] = [x \otimes \phi]. 
	\end{equation}

	Untuk itu, perhatikan bahwa untuk setiap morfisme $f: X \to Y$ dalam
	$\mathcal{A}$, terdapat diagram komutatif
	\[\begin{tikzcd}
		\omega(X) \arrow[r, "{\omega(f)}"] \arrow[d, "\sim"' sloped] & \omega(Y) \arrow[d, "\sim" sloped] \\
		\omega({}^* X)^\vee \arrow[r, "{\omega({}^* f)^\vee}"' inner sep=0.6em] & \omega({}^* Y)^\vee
	\end{tikzcd} \quad \begin{tikzcd}
		\omega(Y)^\vee \arrow[r, "{\omega(f)^\vee}" inner sep=0.6em] \arrow[d, "\sim"' sloped] & \omega(X)^\vee \arrow[d, "\sim" sloped] \\
		\omega({}^* Y) \arrow[r, "{{}^* f}"'] & \omega({}^* X),
	\end{tikzcd}\]
	dengan semua panah vertikal berasal dari isomorfisme dalam paragraf
	sebelumnya. Untuk $\psi \in \omega(Y)^\vee$ dan
	$x \in \omega(X)$, terdapat unsur yang bersesuaian
	$\psi' \in \omega({}^* Y)$ dan $x' \in \omega({}^* X)^\vee$, sedemikian sehingga berturut-turut
	\[ \omega(f)^\vee (\psi) \leftrightarrow \omega({}^* f)(\psi'), \quad \omega(f)(x) \leftrightarrow \omega({}^* f)^\vee(x'). \]
	Dengan memasukkan hal ini ke dalam konstruksi dan deskripsi
	$\coEnd(\omega)$ pada bukti Proposisi
	\ref{prop:cogebra-L-universal}, terlihat bahwa $\delta_f$ di sana
	bersesuaian dengan $\delta_{{}^* f}$. Jadi
	% Koreksi salah ketik pasti pada baris sumber 1578: kelas asal ialah
	% [\phi \otimes x], sesuai dengan definisi S dalam persamaan di atas.
	$[\phi \otimes x] \mapsto [x \otimes \phi]$ memang mendefinisikan secara
	unik pemetaan linear $S$, sehingga \eqref{eqn:Tannaka-aux2}
	terbukti.
	
	Konstruksi pada paragraf sebelumnya juga dapat dilakukan dengan $X^*$
	sebagai ganti ${}^* X$, sehingga memberikan pemetaan linear
	$T: \coEnd(\omega) \to \coEnd(\omega)$; pemetaan ini juga dituliskan
	dalam bentuk $[\phi \otimes x] \mapsto [x \otimes \phi]$. Berdasarkan
	${}^* (X^*) = X = ({}^* X)^*$, mudah dilihat bahwa
	$ST = \identity = TS$; khususnya, $S$ merupakan isomorfisme. Ini adalah
	syarat pertama untuk antipoda.
	
	Selanjutnya, nyatakan perkalian, komultiplikasi, unit, dan homomorfisme
	kounit $\coEnd(\omega)$ masing-masing sebagai $\mu$, $\Delta$, $\eta$,
	$\epsilon$. Pilih objek $X$ dan basis $v_1, \ldots, v_n$ dari ruang
	vektor-$\Bbbk$ $\omega(X)$. Nyatakan basis dual dari $\omega(X)^\vee$
	sebagai $\check{v}_1, \ldots, \check{v}_n$. Nyatakan morfisme-morfisme
	untuk dual kiri $X$ sebagai
	\[ \mathrm{ev}_X: {}^* X \otimes X \to \munit, \quad \mathrm{coev}_X: \munit \to X \otimes {}^* X. \]
	Citra keduanya di bawah $\omega$ merealisasikan dualitas antara
	$\omega(X)$ dan $\omega({}^* X)$. Ingat bahwa $(\cdot)^\vee$ merupakan
	dual dalam kategori monoidal simetris $\cate{Vect}_{\mathrm{f}}(\Bbbk)$,
	tanpa pembedaan kiri dan kanan; mulai sekarang kita dapat
	mengidentifikasikan
	\begin{equation*}
		\omega(\munit) \simeq \omega(\munit)^\vee, \quad \omega(X \otimes {}^* X) \simeq \omega({}^* X \otimes X)^\vee,
	\end{equation*}
	dan, secara bersesuaian, mengidentifikasikan $\omega(\mathrm{coev}_X)$
	dengan $\omega(\mathrm{ev}_X)^\vee$.
	
	Untuk setiap $\phi \in \omega(X)^\vee$ dan $x \in \omega(X)$, deskripsi
	$\Delta$ (atau $\mu$) dalam \eqref{eqn:L-comult} (atau
	\eqref{eqn:L-tensor}) memberikan
	\begin{align*}
		\Delta([\phi \otimes x]) & = \sum_{i=1}^n [\phi \otimes v_i] \otimes [\check{v}_i \otimes x], \\
		\mu(S \otimes \identity)\Delta([\phi \otimes x]) & = \sum_{i=1}^n \mu\left( [v_i \otimes \phi] \otimes [\check{v}_i \otimes x] \right) \\
		& = \sum_{i=1}^n \left[ (v_i \otimes \check{v}_i) \otimes (\phi \otimes x) \right] \\
		& = \left[ \mathrm{coev}_{\omega(X)}(1) \otimes (\phi \otimes x) \right] \\
		& = \left[ \omega(\mathrm{coev}_X)(1) \otimes (\phi \otimes x) \right].
	\end{align*}
	Namun, \eqref{eqn:cogebra-L} menyiratkan bahwa diagram berikut komutatif:
	\[\begin{tikzcd}
		\omega(\munit) \otimes \omega({}^* X \otimes X) \arrow[r, "{\omega(\mathrm{coev}_X) \otimes \identity}" inner sep=0.6em] \arrow[d, "{\identity \otimes \omega(\mathrm{ev}_X)}"'] & \omega(X \otimes {}^* X) \otimes \omega(X \otimes {}^* X), \arrow[d, "{[\cdot]}"] \\
		\omega(\munit) \otimes \omega(\munit) \arrow[d, "\sim"' sloped] \arrow[r, "{[\cdot]}"] & \coEnd(\omega) \\
		\Bbbk \arrow[ru, "\eta"'] & 
	\end{tikzcd}\]
	Oleh karena itu,
	\begin{gather*}
		\left[ \omega(\mathrm{coev}_X)(1) \otimes (\phi \otimes x) \right]
		= \eta\left( \omega(\mathrm{ev}_X)(\phi \otimes x) \right)
		\xlongequal{\text{\eqref{eqn:L-counit}}} \eta\epsilon\left( \phi \otimes x \right),
	\end{gather*}
	yang membuktikan $\mu(S \otimes \identity)\Delta = \eta\epsilon$.
	Dengan cara serupa,
	$\mu(\identity \otimes S)\Delta = \eta\epsilon$. Jadi $S$ merupakan
	antipoda.
\end{proof}

\begin{corollary}\label{prop:Hopf-vs-dual}
	Di bawah bijeksi yang diberikan oleh Teorema
	\ref{prop:Tannaka-reconstruction},
	\[\begin{tikzcd}
		\Obj(\cate{Fib}^{\otimes}(\Bbbk))/\simeq \arrow[shift left, r, "{1:1}"] & \Obj(\Bbbk\dcate{biAlg}) /\simeq , \arrow[shift left, l]
	\end{tikzcd}\]
	$\mathcal{A}$ dalam objek $(\mathcal{A}, \omega)$ di ruas kiri kaku
	kiri dan kaku kanan jika dan hanya jika objek $L$ di ruas kanan merupakan
	aljabar Hopf.
\end{corollary}
\begin{proof}
	Arah ``hanya jika'' mengikuti dari Teorema \ref{prop:Tannaka-aux2},
	sedangkan arah ``jika'' telah dijelaskan dalam Proposisi
	\ref{prop:Hopf-module-dual}.
\end{proof}

Dalam bukti Teorema \ref{prop:Tannaka-aux2}, kekakuan kiri $\mathcal{A}$
memainkan peran utama; kekakuan kanan hanya digunakan untuk memastikan bahwa
$S$ merupakan isomorfisme. Jika definisi aljabar Hopf mengizinkan $S$ yang
tidak dapat dibalik, maka kategori komodul kanan
$\catesubd{Comod}{\mathrm{f}}L$ hanya kaku kiri; lihat rumus dalam Proposisi
\ref{prop:Hopf-module-dual}.
